\chapter{Technical Complements on Positive Maps}
\label{ch:positive_maps_technical_complements}

This appendix collects three groups of technical results used in
Chapter~\ref{ch:channels}.  They establish the inverse-square-root
normalization of positive maps, the two bounds behind Ky Fan's maximum
principle, and the passage from Kraus representations to entrywise complete
positivity.

Throughout, $T^*$ denotes the adjoint for the trace pairing, $\Id$ is the
identity matrix, and $A^{-1/2}$ denotes the inverse of the positive square
root of a positive definite matrix.  The notation $S_k(A)$ is that of
Definition~\ref{def:ky_fan_norm}.

\section{Trace normalization by an inverse square root}

The following results provide the normalization and invertibility steps used
in the proof of Lemma~\ref{lem:making_positive_maps_trace_preserving}.

\begin{lemma}[Trace-normalization criterion, trace-preservation component]
    \label{lem:trace_normalization_criterion_trace}
    \lean{trace_comp_unitaryConjLM_of_conj_traceAdjointMap_one}
    \leanok
    \uses{def:trace_preserving_rectangular, def:trace_pairing_adjoint}
    Let $T:\MN{D}\to\MN{D'}$ be a linear map and let $X\in\MN{D}$ satisfy
    $X^\dagger T^*(\Id)X=\Id$. Then, for every $\rho\in\MN{D}$,
    \begin{align}
        \tr\!(T(X\rho X^\dagger))
        &=\tr(\rho).
        \notag
    \end{align}
\end{lemma}

\begin{proof}\leanok
    Expanding with the trace-pairing adjoint gives
    \begin{align}
        \tr\!(T(X\rho X^\dagger))
        &=\tr\!(X\rho X^\dagger T^*(\Id))
        =\tr\!(\rho X^\dagger T^*(\Id)X)
        =\tr(\rho).
        \notag
    \end{align}
\end{proof}

\begin{theorem}[Filtered trace-normalization criterion]
    \label{thm:filtered_trace_normalization_criterion}
    \lean{comp_unitaryConjLM_positive_tracePreserving_of_conj_traceAdjointMap_one}
    \leanok
    \uses{def:positive_map_rectangular, def:trace_preserving_rectangular,
        def:trace_pairing_adjoint, lem:trace_normalization_criterion_trace,
        thm:conjugation_filters_preserve_positivity}
    Let $T:\MN{D}\to\MN{D'}$ be a positive map between matrix algebras of
    possibly different dimensions, and let $X\in\MN{D}$ satisfy
    $X^\dagger T^*(\Id)X=\Id$, where $T^*$ is the trace-pairing adjoint.
    Then $\rho\mapsto T(X\rho X^\dagger)$ is positive and trace-preserving.

    This is the algebraic trace-normalization step in the proof of
    \cite[Chapter~3, Lemma ``Making positive maps trace preserving'']
    {Wolf2012Quantum}. The inverse-square-root choice is recorded in the
    next theorem.
\end{theorem}

\begin{proof}\leanok
    Positivity follows because $X\rho X^\dagger$ is positive whenever
    $\rho$ is positive, and $T$ is positive. For trace preservation, the
    trace-pairing adjoint identity and cyclicity of the trace give
    \begin{align}
        \tr\!(T(X\rho X^\dagger))
        &=\tr\!(T^*(\Id)X\rho X^\dagger)
        =\tr\!(X^\dagger T^*(\Id)X\rho)
        =\tr(\rho).
        \notag
    \end{align}
\end{proof}

\begin{lemma}[Normalizing conjugation by inverse square root]
    \label{lem:conj_inv_sqrt_normalizes}
    \lean{Matrix.conjTranspose_inv_sqrt_mul_self_mul_inv_sqrt_eq_one_of_posDef}
    \leanok
    \uses{}
    Let $A\in\MN{D}$ be positive definite and set $S=A^{1/2}$. Then
    \begin{align}
        (S^{-1})^\dagger A S^{-1}
        &=\Id.
        \notag
    \end{align}
\end{lemma}

\begin{proof}\leanok
    Since $A$ is positive definite, $S$ is invertible and self-adjoint, and
    $S^2=A$. Therefore
    $(S^{-1})^\dagger A S^{-1}=S^{-1}S^2S^{-1}=\Id$.
\end{proof}

\begin{lemma}[Inverse square root of a positive definite matrix is invertible]
    \label{lem:pos_def_inv_sqrt_invertible}
    \lean{Matrix.PosDef.isUnit_det_inv_sqrt}
    \leanok
    \uses{}
    If $A\in\MN{D}$ is positive definite, then $(A^{1/2})^{-1}$ is
    invertible.
\end{lemma}

\begin{proof}\leanok
    A positive definite matrix is invertible, hence so is its square root
    $A^{1/2}$, and the inverse of an invertible matrix is invertible.
\end{proof}

\begin{theorem}[Trace normalization by inverse square root]
    \label{thm:trace_normalization_inverse_square_root}
    \lean{comp_unitaryConjLM_inv_cfc_sqrt_traceAdjointMap_one}
    \leanok
    \uses{def:positive_map_rectangular, def:trace_preserving_rectangular,
        def:trace_pairing_adjoint,
        thm:filtered_trace_normalization_criterion, lem:conj_inv_sqrt_normalizes}
    Let $T:\MN{D}\to\MN{D'}$ be a positive map such that $T^*(\Id)$ is positive
    definite. Put $X=(T^*(\Id))^{-1/2}$. Then
    $\rho\mapsto T(X\rho X^\dagger)$ is positive and trace-preserving.
\end{theorem}

\begin{proof}\leanok
    Let $S=(T^*(\Id))^{1/2}$. Since $T^*(\Id)$ is positive definite, $S$ is
    invertible and self-adjoint, and $S^2=T^*(\Id)$. For $X=S^{-1}$,
    \begin{align}
        X^\dagger T^*(\Id)X
        &=S^{-1}S^2S^{-1}=\Id.
        \notag
    \end{align}
    Theorem~\ref{thm:filtered_trace_normalization_criterion} applies.
\end{proof}

These results supply the deferred technical steps in
Lemma~\ref{lem:making_positive_maps_trace_preserving}: the inverse-square-root
filter satisfies the normalization identity and is invertible.

\section{Attainment and upper bounds for largest-eigenvalue sums}

The next two results prove attainment and the matching upper bound used in
Theorem~\ref{thm:ky_fan_maximum_principle}.

\begin{theorem}[Achievability of the largest-eigenvalue sum]
    \label{thm:ky_fan_achievability}
    \lean{Matrix.IsHermitian.exists_isProj_trace_eq_kyFanNorm}
    \leanok
    \uses{def:ky_fan_norm}
    For every Hermitian $A\in\MN{D}$ and every $k\ge0$, there is an orthogonal
    projection $P$ of rank $\min(k,D)$ with $\re\tr(PA)=S_k(A)$.
\end{theorem}

\begin{proof}\leanok
    \uses{def:ky_fan_norm}
    Diagonalize $A=U\diag(\lambda_i)U^\dagger$ and take $P$ to be the
    projection onto the span of the eigenvectors with the $\min(k,D)$ largest
    eigenvalues. Then $\re\tr(PA)=\sum_{i\le k}\lambda_i=S_k(A)$.
\end{proof}

\begin{theorem}[Upper bound for the largest-eigenvalue sum]
    \label{thm:ky_fan_upper_bound}
    \lean{Matrix.IsHermitian.trace_mul_re_le_kyFanNorm}
    \leanok
    \uses{def:ky_fan_norm}
    For every Hermitian $A\in\MN{D}$ and every orthogonal projection $P$ of
    rank $k<D$, one has $\re\tr(PA)\le S_k(A)$.
\end{theorem}

\begin{proof}\leanok
    \uses{def:ky_fan_norm}
    Conjugating $P$ by $U$ gives an orthogonal projection $Q$ of the same
    rank, and $\re\tr(PA)=\sum_iw_i\lambda_i$, where $w_i$ are the diagonal
    entries of $Q$. Each $w_i$ lies in $[0,1]$ and the weights sum to
    $\tr(P)=k$. With the threshold $c=\lambda_{k+1}$,
    \begin{align}
        \sum_{i\le k}\lambda_i-\sum_iw_i\lambda_i
        &=\sum_{i\le k}(\lambda_i-c)(1-w_i)
          +\sum_{i>k}(c-\lambda_i)w_i
        \ge0,
        \notag
    \end{align}
    since for $i\le k$ the eigenvalues dominate $c$ while $1-w_i\ge0$, and
    for $i>k$ they are dominated by $c$ while $w_i\ge0$.
\end{proof}

Together, Theorems~\ref{thm:ky_fan_achievability} and
\ref{thm:ky_fan_upper_bound} give the two inequalities in
Theorem~\ref{thm:ky_fan_maximum_principle}.

\section{From Kraus representations to complete positivity}

These results connect the Kraus formulation in
Definition~\ref{def:cp_map} with entrywise positivity on block matrices and
with the corresponding completely positive map between matrix
$C^*$-algebras.

\begin{theorem}[Entrywise complete positivity from a Kraus representation]
    \label{thm:cp_map_entrywise_complete_positive}
    \lean{IsCPMap.map_cstarMatrix_nonneg}
    \leanok
    \uses{def:cp_map}
    Let $E : \MN{D} \to \MN{D}$ be completely positive in the Kraus sense.
    For every $k$ and every positive block matrix
    $M \in M_k(\MN{D})$, the entrywise image
    $(E(M_{ab}))_{a,b}$ is again positive in $M_k(\MN{D})$.
\end{theorem}

\begin{proof}\leanok
    Write $E$ in the Kraus form (\ref{eq:channel_kraus_representation}).
    For each $i$, let $d_i$ be the block-diagonal matrix with $K_i$ on every
    diagonal block. Then
    \begin{align}
        (E(M_{ab}))_{a,b}
        &=\sum_i d_i M d_i^\dagger.
        \notag
    \end{align}
    Each term on the right is positive because conjugation preserves
    positivity, and a finite sum of positive matrices is positive.
\end{proof}

\begin{theorem}[Kraus maps as abstract completely positive maps]
    \label{thm:cp_map_to_abstract_completely_positive}
    \lean{IsCPMap.toCompletelyPositiveMap}
    \leanok
    \uses{def:cp_map, thm:cp_map_entrywise_complete_positive}
    Every Kraus-represented completely positive map
    $E : \MN{D} \to \MN{D}$ determines a completely positive map between
    the matrix $C^*$-algebras in the entrywise positivity sense.
\end{theorem}

\begin{proof}\leanok
    The defining entrywise positivity condition is exactly
    Theorem~\ref{thm:cp_map_entrywise_complete_positive}.
\end{proof}

Thus the Kraus formulation of Definition~\ref{def:cp_map} has the two
consequences cited there: entrywise complete positivity and the associated
abstract completely positive map.
